\par\indent\SourceLeaf{209}%
Es wird daher, wenn wir wiederum durch die Charakteristik \(f\) eine willkürliche Funktion andeuten, \(X\) dem reellen und \(iY\) dem imaginären Theile von
\[
f(t+i\log\mathrm{cotang}\tfrac{1}{2}u)
\]
\noindent gleich gesetzt werden müssen. Wir wollen ein Paar specielle Fälle dieser allgemeinen Auflösung anführen.

Wählt man für \(f\) eine lineäre Function, indem man \(f\upsilon=k\upsilon\) setzt, so wird
\[
X=kt,\quad Y=k\log\mathrm{cotang}\tfrac{1}{2}u
\]
\noindent Auf die Erde angewandt, ist dies, wenn man \(t\) die geographische Länge, \(90^0-u\) die Breite bedeuten lässt, offenbar mit \textsc{Mercators} Projection einerlei. Für das Vergrösserungsverhältniss geben hier die Formeln des 7. Artikels
\[
m=\frac{k}{a\sin u}
\]

Nimmt man für \(f\) eine imaginäre Exponentialfunction, und zwar zuerst die einfachste \(f\upsilon=ke^{i\upsilon}\), so wird
\[
f(t+i\log\mathrm{cotang}\tfrac{1}{2}u)=ke^{\log\mathrm{tang}\frac{1}{2}u+it}=k\,\mathrm{tang}\tfrac{1}{2}u(\cos t+i\sin t)
\]
\noindent und
\[
X=k\,\mathrm{tang}\tfrac{1}{2}u\mathord{.}\cos t,\quad
Y=k\,\mathrm{tang}\tfrac{1}{2}u\mathord{.}\sin t
\]
\noindent welches, wie man leicht sieht, die stereographische Polarprojection ist.

Setzt man allgemeiner \(f\upsilon=ke^{i\lambda\upsilon}\), so wird
\[
X=k\,\mathrm{tang}\tfrac{1}{2}u^{\lambda}\mathord{.}\cos\lambda t,\quad
Y=k\,\mathrm{tang}\tfrac{1}{2}u^{\lambda}\mathord{.}\sin\lambda t
\]
\noindent Für das Vergrösserungsverhältniss erhalten wir hier
\[
n=aa\sin u^2,\quad N=1,\quad\varphi\upsilon=i\lambda ke^{i\lambda\upsilon}
\]
\noindent und hieraus
\[
m=\frac{\lambda k\,\mathrm{tang}\frac{1}{2}u^{\lambda}}{a\sin u}
\]

Man sieht, dass hier die Darstellung aller Punkte, für welche \(u\) constant ist, in Einen Kreis, und die Darstellung aller Punkte, für welche \(t\) constant ist, in Eine gerade Linie fällt, wie auch, dass die allen verschiedenen Werthen von \(u\) angehörigen Kreise concentrisch sind. Dies gibt eine sehr zweckmässige Kartenprojection, wenn nur ein Theil der Kugelfläche darzustellen ist, und man thut dann am besten, \(\lambda\) so zu wählen, dass das Vergrösserungsverhältniss für die
